\documentclass[10pt,a4paper]{amsart}
\usepackage[margin=19mm]{geometry}
\usepackage{amsmath,amssymb,mathtools}
\usepackage[scr=rsfs,cal=euler]{mathalfa}
\usepackage{xfrac}
\usepackage[unicode,hidelinks,bookmarks=false]{hyperref}
\newcommand{\ie}{i.e.\ }
\newcommand{\cf}{cf.\ }
\newcommand{\resp}{respectively\ }
\newcommand{\Subsection}{\subsection}
\providecommand{\nobreakdash}{\nobreak\hbox{-}}
\setlength{\emergencystretch}{2em}
\multlinegap0pt
\usepackage{bm}
\usepackage{bbm}
\usepackage{dsfont}
\usepackage{xspace}
\usepackage{cancel}
\usepackage{mathtools}
\usepackage{amsthm}
\makeatletter
\@ifundefined{theorem}{\newtheorem{theorem}{Theorem}}{}
\makeatother
\newcommand{\PSL}{\mathrm{PSL}}
\newcommand{\SL}{\mathrm{SL}}
\newcommand{\calN}{\mathcal{N}}
\newcommand{\calQ}{\mathcal{Q}}
\newcommand{\calR}{\mathcal{R}}
\newcommand{\calX}{\mathcal{X}}
\title[Computations of parabolic character schemes of knots]{Computations of parabolic character schemes of knots}
\author{Yunhi Cho, Hyuk Kim, Seonhwa Kim, Seokbeom Yoon}
\date{Source: arXiv:2608.06890v1 (CC BY 4.0)}
\begin{document}
\maketitle

\noindent\emph{Source and licence.} Computations of parabolic character schemes of knots. Version arXiv:2608.06890v1 (CC BY 4.0), distributed under the Creative Commons Attribution 4.0 International licence, as recorded in the arXiv licence metadata for this exact version and shown on the abstract page https://arxiv.org/abs/2608.06890v1. Full rights notice: licenses/arxiv-2608.06890-CC-BY-4.0.txt.\par\smallskip

\noindent\textit{Editorial scope.} Counted unit: the Theorem whose source label is \texttt{thm:normalizedchart} in the article (source0.tex lines 859--870, source label \texttt{thm:normalizedchart}), delivered together with the article's own complete original proof of it (source0.tex lines 872--923, one \texttt{proof} environment of 52 lines). No mathematics is written, repaired or re-derived: statement and proof are the article's own text line for line. Required context retained uncounted: prop.mu2torsor (lines 751--754). Every definition, display and label of those lines is retained unchanged. Omitted from this package and not redistributed: 1--750, 755--858, 871--871, 924--1381. Nothing inside a retained excerpt is omitted or altered: every retained body is delivered exactly as the article's lines are written, so no omission receipt of any kind accompanies this package. Citations: the retained excerpts carry no citation command at all, so no bibliography entry of the article is transcribed and no reference is redistributed. Reference adaptation: none was needed and none was made; every reference inside the delivered unit resolves inside it, so no unresolved reference or citation is produced. Printed slips kept literal: no printed word, symbol, hypothesis or number of the counted statement or of its proof was altered. Layout and class: the article's own class is \texttt{amsart} with packages including fullpage, amssymb, amsmath, amscd, bbm, tikz-cd, xy, url, hyperref, slashed, listings, verbatim, mathtools, ulem; the wrapper is the lane's pinned \texttt{amsart} editorial wrapper at 10pt on A4, which restates the theorem-like environments the retained text needs and adds the article's own macro definitions that the retained text needs (\texttt{\textbackslash PSL}, \texttt{\textbackslash SL}, \texttt{\textbackslash calN}, \texttt{\textbackslash calQ}, \texttt{\textbackslash calR}, \texttt{\textbackslash calX}), with extra wrapper packages (bm, bbm, dsfont, xspace, cancel, mathtools). The delivered numbering is the unit's own. Assets: no retained span contains a figure environment, a graphics-inclusion command, a PDF-page inclusion, a table or a TikZ picture; no image file is copied into this package, the package declares no asset dependency and no image file is redistributed with it. Licence: version arXiv:2608.06890v1 of this article is distributed under the Creative Commons Attribution 4.0 International licence, as recorded in the arXiv licence metadata for this exact version and shown on the abstract page https://arxiv.org/abs/2608.06890v1. A full rights notice is delivered at licenses/arxiv-2608.06890-CC-BY-4.0.txt. Characters: every retained excerpt is pure ASCII, so the delivered engine, XeLaTeX with its default Latin Modern text font, reports no missing character.
	\begin{theorem} \label{prop.mu2torsor}
		The morphism $	\rho_{\boldsymbol{\epsilon}}:\calQ_{\boldsymbol{\epsilon}}^\times(D) \rightarrow
		\calR_s^\times(K)$ is a surjective $\mu_2$-torsor.
	\end{theorem}

	\begin{theorem} \label{thm:normalizedchart}
		For $k=2,\ldots,n$, 
		$$
		\Theta_k: \SL_2 \times\calN_{\boldsymbol{\epsilon},k}(D)\longrightarrow \calQ^\times_{{\boldsymbol{\epsilon}},k}(D),\qquad
		(g,\boldsymbol{\alpha})\longmapsto g\boldsymbol{\alpha},
		$$
		is an isomorphism.  % Here $\SL_2 = \Spec ( \BC[a,b,c,d]/(ad-bc-1))$.
		In addition, it induces an isomorphism
		$$
		\calN_{\boldsymbol{\epsilon},k}(D) \overset{\sim}{\longrightarrow} \calX_{s,k}^\times(K)\,.
		$$
	\end{theorem}

	\begin{proof}
		Let	$\boldsymbol{\alpha}\in \calQ^\times_{{\boldsymbol{\epsilon}},k}(D)$ with $
		\alpha_{1}	= \binom{x_1}{y_1}$, $\alpha_{k} =\binom{x_{k}}{y_{k}}$.
		Since $\Delta_k$ is invertible on $\calQ^\times_{{\boldsymbol{\epsilon}},k}(D)$,
		$$
		F(\boldsymbol{\alpha}) 	=
		\begin{pmatrix}
			y_{k}\Delta_k^{-1}	& 	-x_{k}\Delta_k^{-1} \\
			-y_{1}		& 		x_{1}
		\end{pmatrix}
		$$
		defines a regular $\SL_2$-valued function $F$ on $\calQ^\times_{{\boldsymbol{\epsilon}},k}(D)$.
		Then we have
		$$
		F(\boldsymbol{\alpha})\alpha_{1} = \binom{1}{0},	\qquad
		F(\boldsymbol{\alpha})\alpha_{k} = \binom{0}{\Delta_k}.
		$$
		It follows that	$F(\boldsymbol{\alpha})\boldsymbol{\alpha} \in	\calN_{\boldsymbol{\epsilon},k}(D)$. The inverse of $\Theta_k$ is then  given by
		$$
		\boldsymbol{\alpha}	\longmapsto	(F(\boldsymbol{\alpha})^{-1}, \, F(\boldsymbol{\alpha})\boldsymbol{\alpha}).
		$$
		Indeed, if $g\in\SL_2$ satisfies $g\alpha_{1} = \binom{1}{0}$ and $g\alpha_{k} = \binom{0}{c}$ for some $c$,
		then 
		$$
		c =	\det \left(g\alpha_{1},	g\alpha_{k}	\right) = \det \left(\alpha_{1}, \alpha_{k} \right) = \Delta_k.
		$$
		Since $\alpha_1$ and $\alpha_k$ form a basis of
		$\mathcal O({\calQ^\times_{{\boldsymbol{\epsilon}},k}(D)})^{\oplus2 }$,  these conditions determine $g$ uniquely. This proves that $\Theta_k$ is an isomorphism.

		By Proposition~\ref{prop.mu2torsor}, the restriction of $\rho_{\boldsymbol{\epsilon}}: \calQ^\times_{\boldsymbol{\epsilon}}(D) \rightarrow \calR^\times_{s}(K)$ to $\calQ^\times_{{\boldsymbol{\epsilon}},k}(D)$
		$$
		\rho_{\boldsymbol{\epsilon}}: 	\calQ^\times_{{\boldsymbol{\epsilon}},k}(D)\longrightarrow	\calR^\times_{s,k}(K)
		$$
		is a $\mu_2$-torsor. Therefore,
		$\Gamma (\calR^\times_{s,k}(K)) \cong \Gamma (\calQ^\times_{{\boldsymbol{\epsilon}},k}(D))^{\mu_2}$.
		On the other hand,  $\mu_2$ is the center of $\SL_2$ and acts trivially on
		$\calR^\times_{s,k}(K)$. Hence, the $\SL_2$-action factors through
		$\PSL_2=\SL_2/\mu_2$. It follows that
		\[
		\Gamma (\calX_{s,k}^\times(K)) \cong
		\Gamma (\calR^\times_{s,k}(K))^{\PSL_2}	\cong 
        (\Gamma (\calQ^\times_{{\boldsymbol{\epsilon}},k}(D))^{\mu_2})^{\PSL_2} \cong 
        \Gamma (\calQ^\times_{{\boldsymbol{\epsilon}},k}(D))^{\SL_2}.
		\]
		Under the isomorphism $\Theta_k$, the $\SL_2$-action on $\SL_2 \times\calN_{\boldsymbol{\epsilon},k}(D)$ is given by the left multiplication on the first factor. It follows that
		$$
		\Gamma (\calQ^\times_{{\boldsymbol{\epsilon}},k}(D))^{\SL_2} \cong
        \Gamma(\SL_2 \times\calN_{\boldsymbol{\epsilon},k}(D))^{\SL_2} \cong
		\Gamma (\calN_{\boldsymbol{\epsilon},k}(D)).
		$$
		Therefore, we obtain $\calN_{\boldsymbol{\epsilon},k}(D) \cong \calX_{s,k}^\times(K)$.
	\end{proof}

\end{document}
